\chapter{Literature Survey}

\section{Definitions and Importance of Churn}

Customer churn, also referred to as customer attrition or customer defection, is broadly defined
as the process by which customers cease their relationship with a business entity over a given
observation period. In the FinTech domain, churn can manifest as account closure, cessation of
transactional activity, or migration to a competing platform~\cite{ngai2009application}. The
distinction between \textit{voluntary churn} (customer-initiated departure due to dissatisfaction
or superior alternatives) and \textit{involuntary churn} (service termination due to payment
failure or fraud) is important from a modelling perspective, as the predictive signals differ
substantially between the two categories.

The economic importance of churn prevention has been extensively quantified in the literature.
Gupta et al.~\cite{gupta2006modeling} demonstrated that a one-percentage-point reduction in churn
rate generates a greater improvement in firm value than a comparable reduction in customer
acquisition cost or variable cost, underscoring the strategic priority of retention. In the
financial services sector, the situation is particularly acute: the average churn rate for digital
banking applications exceeds twenty-five percent annually, and the average cost of replacing a
churned customer, when factoring in marketing expenditure, onboarding overhead, and the revenue
forgone during the transition period, is estimated at between four hundred and twelve hundred
US dollars~\cite{hung2006applying}.

Early churn prediction research relied heavily on simple threshold-based rules applied to a small
number of manually selected indicators, such as recency of last purchase or frequency of logins.
While interpretable, these approaches suffer from low recall---they miss a significant proportion
of churners who do not exhibit the specific thresholded behaviours---and from high false-positive
rates that result in unnecessary campaign expenditure. The subsequent shift to machine-learning
methods has enabled substantially richer feature utilisation and more accurate probabilistic risk
scoring.

\section{Traditional Machine Learning Models}

\subsection{Logistic Regression}

Logistic Regression (LR) was one of the earliest ML techniques applied to churn
prediction~\cite{burez2009handling}. Its principal advantage is interpretability: the coefficient
of each predictor directly quantifies the direction and magnitude of its association with churn
risk on the log-odds scale. Coussement and Van den Poel~\cite{coussement2008churn} applied LR to
a Belgian newspaper subscription dataset and reported an area under the ROC curve (AUC) of 0.81,
establishing a meaningful baseline. However, LR assumes a linear relationship between features
and the log-odds of churn, a restriction that limits its capacity to capture the complex
non-linear interactions present in FinTech data.

\subsection{Decision Trees}

Decision Tree classifiers partition the feature space into axis-aligned rectangular regions through
a recursive greedy splitting procedure, typically optimised by minimising Gini impurity or
information gain. Their primary virtue is interpretability: the resulting tree structure can be
visualised and communicated to business stakeholders without requiring statistical expertise.
Lemmens and Croux~\cite{lemmens2006bagging} evaluated decision trees for churn prediction in
the telecommunications sector and noted a tendency towards overfitting when trees are grown to
excessive depth. Regularisation through minimum leaf size constraints and maximum depth limits
partially mitigates this, but single-tree classifiers remain outperformed by ensemble methods.

\subsection{Support Vector Machines}

Support Vector Machines (SVM) seek to identify the maximum-margin hyperplane separating positive
(churn) and negative (retain) instances in the feature space, potentially after applying a kernel
function to map inputs into a higher-dimensional space where linear separation is feasible. Huang
et al.~\cite{huang2015customer} applied radial-basis-function SVMs to a telecommunications churn
dataset and achieved AUC values competitive with ensemble methods. The main practical limitation
of SVMs in the churn prediction context is their computational cost during both training and
prediction, which scales unfavourably with dataset size.

\section{Ensemble and Boosting Methods}

\subsection{Random Forest}

Breiman's Random Forest~\cite{breiman2001random} combines the predictions of many de-correlated
decision trees grown on bootstrap samples of the training data using random feature subsets at
each split. The aggregation of diverse trees substantially reduces variance without a commensurate
increase in bias, yielding models that generalise well to unseen data. Verbeke et al.~\cite{verbeke2012new}
demonstrated that Random Forest achieves superior top-decile lift---a metric of particular
commercial relevance in targeted retention campaigns---compared to both single-tree and LR
classifiers.

\subsection{Gradient Boosting Machines}

Gradient Boosting Machines (GBM), introduced by Friedman~\cite{friedman2001greedy}, construct an
additive ensemble of shallow trees by iteratively fitting each new tree to the pseudo-residuals
(negative gradients of the loss function) of the current ensemble. Gradient boosting is less
susceptible to overfitting than simple decision trees and achieves higher predictive accuracy than
Random Forest on many structured datasets, at the cost of increased training time and a larger
hyperparameter search space.

\subsection{XGBoost}

Extreme Gradient Boosting (XGBoost), developed by Chen and Guestrin~\cite{xgboost2016}, extends
the standard gradient boosting framework with a regularised objective function that penalises tree
complexity, a novel sparsity-aware split-finding algorithm that handles missing values natively,
and highly optimised parallel computation. XGBoost has dominated Kaggle competitions involving
structured data and has been widely adopted for churn prediction in commercial settings. Stripling
et al.~\cite{stripling2018profit} reported XGBoost as the top-performing algorithm in a benchmark
study covering eight churn datasets, outperforming Random Forest in terms of both AUC and
profit-driven metrics.

\subsection{LightGBM}

LightGBM, developed at Microsoft Research~\cite{ke2017lightgbm}, further accelerates gradient
boosting via two novel techniques: Gradient-based One-Side Sampling (GOSS), which preferentially
retains instances with large gradient magnitudes, and Exclusive Feature Bundling (EFB), which
compresses mutually exclusive sparse features to reduce dimensionality. These innovations enable
LightGBM to train five to twenty times faster than XGBoost on large datasets with minimal loss in
predictive accuracy, making it particularly suitable for real-time churn scoring applications.

\section{Neural Networks and Deep Learning}

\subsection{Multi-Layer Perceptrons}

Multi-Layer Perceptrons (MLP) are feedforward neural networks with one or more hidden layers of
densely connected neurons. When applied to tabular data for churn prediction, MLPs require
careful regularisation (dropout, weight decay) and hyperparameter tuning to avoid overfitting,
particularly when the training dataset is moderately sized. Mozer et al.~\cite{mozer2000predicting}
pioneered the use of neural networks for customer churn prediction in the telecommunications
sector, demonstrating that a three-layer MLP outperformed conventional statistical models on a
proprietary subscriber dataset.

\subsection{Recurrent Neural Networks}

Recurrent Neural Networks (RNN) and their Long Short-Term Memory (LSTM) variant are naturally
suited to sequential data, enabling the model to capture temporal dependencies in transaction
histories that are invisible to batch-aggregation approaches. Wangperawong et al.~\cite{wangperawong2016churn}
modelled individual transaction sequences as time series inputs to an LSTM and demonstrated
improved churn recall relative to static-feature baselines. The computational cost of training
sequence models and the requirement for aligned temporal records are practical constraints in
production deployments.

\subsection{Attention Mechanisms and Transformers}

Recent work has explored Transformer-based architectures~\cite{vaswani2017attention} for tabular
data. TabNet~\cite{arik2021tabnet} uses sequential attention to select a sparse subset of features
at each decision step, yielding simultaneously high predictive performance and instance-level
interpretability. Preliminary evaluations suggest TabNet is competitive with XGBoost on large
FinTech datasets, while its native feature selection mechanism reduces the dependency on manual
feature engineering.

\section{Handling Class Imbalance}

Class imbalance is a pervasive challenge in churn prediction: in most real-world deployments,
churners represent fewer than twenty percent of the active customer base, and in some contexts the
ratio reaches one-to-twenty or more. Standard classifiers trained on imbalanced data are biased
towards the majority class, achieving high accuracy by predominantly predicting ``no churn'' while
delivering low recall for the minority churn class.

\subsection{Resampling Strategies}

Chawla et al.~\cite{chawla2002smote} introduced the Synthetic Minority Over-sampling Technique
(SMOTE), which generates synthetic minority-class instances by interpolating along line segments
connecting existing minority-class nearest neighbours in feature space. SMOTE is widely adopted
and has been shown to improve minority-class recall substantially without simply duplicating
existing examples. Tomek links---pairs of instances from opposite classes that are mutual nearest
neighbours---can be removed as a complementary under-sampling step to clean ambiguous boundary
regions.

\subsection{Cost-Sensitive Learning}

An alternative to resampling is to incorporate class weights directly into the model's loss
function, penalising misclassification of minority-class instances more heavily. Most scikit-learn
classifiers support a \texttt{class\_weight} parameter that can be set to ``balanced'' to
automatically compute weights inversely proportional to class frequencies. Bahnsen et al.~\cite{bahnsen2016feature}
extended cost-sensitive learning by incorporating instance-specific misclassification costs
derived from customer lifetime value estimates, a natural fit for FinTech applications.

\section{Feature Engineering Strategies}

Feature engineering---the process of constructing informative input representations from raw
data---has been consistently identified as one of the most impactful levers for improving churn
model performance~\cite{provost2013data}. Key strategies include:

\begin{itemize}
    \item \textbf{RFM features}: Recency (days since last transaction), Frequency (number of
          transactions in a window), and Monetary value (total spend) provide a compact
          behavioural summary that has proven predictive across diverse churn
          contexts~\cite{fader2005rfm}.

    \item \textbf{Rolling-window aggregation}: Computing statistics (mean, standard deviation,
          trend slope) over multiple time windows (30, 60, 90 days) captures temporal dynamics
          and enables the model to distinguish customers exhibiting declining engagement from
          stable ones.

    \item \textbf{Ratio and interaction features}: Derived quantities such as the credit
          utilisation ratio (outstanding balance divided by credit limit) or the support contact
          rate (tickets per month of tenure) often carry more predictive signal than their
          constituent raw features.

    \item \textbf{Categorical embeddings}: For high-cardinality categorical variables such as
          product category or geographic region, target-encoding (replacing the category with
          the within-category mean of the target variable) or learned entity embeddings may
          outperform one-hot encoding.
\end{itemize}

\section{Evaluation Metrics}

The choice of evaluation metric has a decisive influence on which model is selected for
deployment. Accuracy is inappropriate for imbalanced churn datasets because a trivial model that
always predicts ``no churn'' achieves high accuracy while providing zero actionable value.
The metrics most commonly used in churn prediction literature are:

\begin{itemize}
    \item \textbf{ROC-AUC}: The area under the Receiver Operating Characteristic curve measures
          the probability that the model ranks a randomly chosen churner above a randomly chosen
          non-churner. It is threshold-independent and robust to class imbalance.

    \item \textbf{Precision, Recall, and F1-score}: Precision is the proportion of predicted
          churners that are true churners; Recall (Sensitivity) is the proportion of actual
          churners that are correctly identified. The F1-score is their harmonic mean, providing
          a single balanced metric particularly useful when both false positives (campaign waste)
          and false negatives (missed churners) carry significant cost.

    \item \textbf{Top-decile lift}: The ratio of the churn rate within the top ten percent of
          risk-scored customers to the overall churn rate in the population. It directly quantifies
          the efficiency gain achieved by targeting the highest-risk cohort.

    \item \textbf{Kolmogorov--Smirnov (KS) statistic}: The maximum difference between the
          cumulative distribution of churn probabilities for churners and non-churners. KS
          statistics above 40 are generally considered commercially useful.
\end{itemize}

\section{Explainability and Interpretable AI}

Regulatory frameworks such as the European Union's General Data Protection Regulation (GDPR) and
the Algorithmic Accountability Act impose requirements for explanations of automated decisions
that affect individuals. In the FinTech domain, model explainability is also a prerequisite for
adoption by risk management teams and business stakeholders accustomed to transparent scorecard
models.

SHAP (SHapley Additive exPlanations)~\cite{lundberg2017unified} provides a theoretically grounded
attribution framework based on cooperative game theory. Each feature is assigned a Shapley value
that represents its average marginal contribution to the model's prediction across all possible
feature coalitions. SHAP values satisfy several desirable axioms (local accuracy, missingness,
consistency) that other attribution methods do not. For tree-based models, TreeSHAP~\cite{lundberg2020local}
computes exact Shapley values in polynomial time, making it practical for production-scale
datasets.

LIME (Local Interpretable Model-agnostic Explanations)~\cite{ribeiro2016should} is an alternative
that approximates the complex model locally with a simple interpretable surrogate (e.g., a sparse
linear model) fitted to perturbations of a specific input instance. While LIME is
model-agnostic---applicable to any black-box classifier---it produces explanations that are
sensitive to the perturbation strategy and may be inconsistent across runs.

\section{Summary}

The literature survey reveals several consistent themes:

\begin{enumerate}
    \item Ensemble methods---particularly gradient boosting variants (XGBoost, LightGBM)---achieve
          the highest predictive accuracy on structured FinTech churn datasets.
    \item Class imbalance must be explicitly addressed, with SMOTE-based resampling offering a
          reliable and widely validated approach.
    \item Feature engineering, especially the construction of RFM metrics and rolling-window
          statistics, provides substantial predictive value over raw input features.
    \item Explainability via SHAP values has emerged as the industry standard for communicating
          model decisions to business and regulatory stakeholders.
    \item There is a growing trend towards real-time scoring architectures and streaming feature
          pipelines, which the present work addresses in its deployment design.
\end{enumerate}

The gaps identified in the literature---specifically the lack of comprehensive production-grade
implementations that integrate preprocessing, imbalance handling, interpretable modelling, and
API deployment in a single end-to-end pipeline for the FinTech context---directly motivate the
proposed system described in the subsequent chapters.
